% FlashTeX demo: only supported commands, with UTF-8 text and multiple pages.
\begin{document}
\section{Résumé of a small journey}
This sample follows a notebook through a quiet day of reading, walking, and writing. It begins at a café beside the station, where a traveller opens a fresh page and records the morning light. The aim is simple: give the preview enough natural prose to wrap across many lines, carry words onto another page, and keep every visible word connected to its original place in the source.

The first entry makes a \textbf{clear observation}: small details reward patient attention. Steam rises from a cup, a bicycle leans against the window, and the door opens with a soft bell. A naïve sketch of the scene becomes more useful when the writer adds a few careful sentences. There is no need for elaborate notation here — ordinary words already offer plenty to read and inspect.

\subsection{At the café}
Across the room, two readers compare their notes about a familiar story. One remembers the setting while the other remembers a conversation. Their different accounts suggest that a useful notebook leaves room for several views. The traveller writes an \emph{important reminder} in the margin of the entry: describe what happened before deciding what it means, and return later if the first explanation feels too easy.

Outside, the street grows busy as shops lift their shutters and deliveries arrive. A baker carries warm bread past the station steps while a neighbour pauses to ask about the weather. The notebook gathers these small events into one paragraph, then gives the next thought its own space. Blank lines in the source make those pauses easy to see when reading or editing the document.

\textit{A note before leaving}: the café is a useful place to begin, but the rest of the town deserves a closer look. The traveller closes the notebook, checks the route, and walks toward the river. This sentence ends with an explicit paragraph command.\par
The next paragraph begins here without a blank line in the source. It describes a bridge with a broad footpath, low stone walls, and a view toward the old library. Water reflects the pale sky beneath it. A slow walk gives the writer time to choose precise words and notice how a single scene can support several different descriptions without becoming a list of disconnected facts.

\section{Pages from the river walk}
At the bridge, the traveller stops to watch leaves move along the water. Some turn in small circles near the bank; others pass steadily beneath the arch. The notebook compares these movements with the pace of a working day. A short pause can clarify a difficult thought, while a longer stretch of steady effort can bring an unfinished idea closer to a form that another person can understand.

The path continues beside a row of trees, each casting a different pattern on the pavement. Children count the ducks near a landing, and a gardener gathers fallen branches into a cart. These observations provide enough continuous text for the compiler to break lines naturally. Each paragraph also offers recognizable phrases, so a reader can select a word in the preview and find the corresponding passage in the editor.

Two short lines mark a pause in the walk.\\
The river carries the morning onward.\\
The notebook carries a few memories home.

\subsection{Reading at the library}
Inside the library, a long table stands beneath tall windows. The traveller chooses a seat near the end and reads an account of the town from many years ago. Familiar streets appear under older names, and a busy square was once an orchard. The contrast encourages another careful paragraph about change: places retain traces of earlier lives even when the people who remember those lives have moved away.

A librarian brings a second volume with a simple map drawn inside the cover. The traveller studies it for a while and then records the route in plain prose. The description begins at the station, passes the café, crosses the bridge, and ends beside the reading room. Repeating a few landmarks gives the document a thread that is easy to follow across a page boundary without needing special layout commands.

The afternoon entry returns to the words naïve and Résumé, along with café and an em dash — all written directly as Unicode text. These words are part of the story as well as useful examples for checking source positions. A reader should be able to move from an accented word in the preview to the same word in the editor, even when earlier characters occupy more than one byte.

\section{Bringing the notebook home}
By late afternoon, the traveller has gathered enough material for a short account of the day. The first draft is deliberately modest. It records places, conversations, and a few changes of mind without trying to turn every observation into a lesson. A \textbf{useful revision} will keep the details that help a reader picture the journey and shorten the sentences that merely repeat what the previous paragraph already explained.

At home, the notebook rests beside a lamp while the traveller reads the pages once more. The opening scene now feels connected to the river and the quiet library. A friend could follow the route from these paragraphs alone, imagining the sound of the bell, the leaves on the water, and the light across the reading table. That continuity is the small achievement the writer hoped to preserve.

The final entry leaves space for tomorrow. Another walk may reveal a different street, another book may change an assumption, and another conversation may offer a better question. For now, the document ends with enough prose to exercise wrapping and pagination, clear headings to guide the reader, and ordinary paragraphs that can be edited and compiled again. The journey is complete, but the notebook remains ready for its next page.
\end{document}
